\chapter{Conclusiones y líneas futuras}
\label{ch:conclusiones}

%--------------------------------------------------------------------
\section{Conclusiones}
%--------------------------------------------------------------------

Este trabajo ha abordado el problema de la optimización de carteras según
el modelo de Markowitz, prestando especial atención a la estimación de sus
dos parámetros fundamentales: la matriz de covarianza $\boldsymbol{\Sigma}$
y el vector de rentabilidades esperadas $\boldsymbol{\mu}$. A continuación
se sintetizan las principales conclusiones.

\subsection{Sobre la formulación matemática}

El modelo de Markowitz es un programa cuadrático convexo (QP) de solución exacta
cuyas condiciones KKT admiten una lectura económica clara: en el óptimo, el
riesgo adicional que aporta cada activo es proporcional a su rentabilidad
esperada. La excepción son los activos que chocan con una de las cotas de peso
(quedarse en el $0$ o en el tope del $20\%$); para ellos, el precio sombra de
esa cota mide cuánto bajaría el riesgo de la cartera si la cota pudiera
relajarse. Sin restricciones de desigualdad, la frontera eficiente es una hipérbola
en el plano $(\sigma_p, \mu_p)$ cuyos parámetros dependen de tres escalares
($A$, $B$, $C$) que condensan la información de $\boldsymbol{\mu}$ y
$\boldsymbol{\Sigma}$.

\subsection{Sobre la estimación de la covarianza}

La matriz de covarianza muestral, aunque insesgada, produce
carteras excesivamente concentradas e inestables cuando el número de activos
es comparable al de observaciones ($n/T \approx 0{,}83$ en nuestro caso).
Se han evaluado dos alternativas. El \textbf{shrinkage de Ledoit-Wolf} reduce el
error cuadrático medio y da carteras más diversificadas y con mejor
comportamiento out-of-sample (Sharpe $0{,}68$ frente a $0{,}64$ del muestral). El
\textbf{estimador PCA}, que retiene los factores principales y conserva el riesgo
específico de cada activo, obtiene el mejor Sharpe de las estrategias clásicas
($0{,}69$), aunque más concentrado.

La equivalencia entre el estimador PCA y el modelo factorial APT de Ross
\cite{Ross1976} constituye uno de los resultados metodológicos más
relevantes del trabajo: los componentes principales de la matriz de
covarianza son, bajo condiciones adecuadas, los factores de riesgo
sistemático del mercado.

\subsection{Sobre la predicción de \texorpdfstring{$\boldsymbol{\mu}$}{μ} con Machine Learning}

Se ha implementado un sistema de predicción de $\boldsymbol{\mu}$ con tres
modelos (Ridge, Lasso y Random Forest) entrenados de forma agrupada sobre los
30 activos, con 11 features derivadas exclusivamente de precios.

\textbf{Lasso} es el mejor modelo ML: su Sharpe de $0{,}75$ iguala en la
práctica a la $1/N$, con el menor turnover (la rotación de la cartera de un mes
a otro; $0{,}12$) y la mayor diversificación de las carteras optimizadas; la
regularización $\ell_1$ es especialmente valiosa en entornos de
baja señal-ruido. \textbf{Random Forest} ($0{,}71$) y \textbf{Ridge} ($0{,}68$)
son también competitivos; Ridge logra además el menor drawdown del estudio
($18{,}4\%$). Ninguna de estas
diferencias, tampoco frente a los baselines, es estadísticamente significativa:
con 143 meses, los intervalos de confianza del Sharpe son demasiado anchos para
distinguir las estrategias.

La conclusión más relevante en este ámbito es que \textbf{la viabilidad del ML
depende sobre todo de compartir información entre activos}: al estimar un
único modelo agrupado (en torno a 1.000 observaciones, en lugar de 36 por
activo), incluso un método no lineal como Random Forest se estabiliza y deja
de sobreajustar.

\subsection{Sobre la comparación con baselines pasivos}

El resultado más honesto de este trabajo, y quizá el más valioso, es que
\textbf{ninguna cartera optimizada supera de forma estadísticamente
significativa a los baselines pasivos}: aunque Lasso iguala a la $1/N$ (Sharpe
$0{,}75$ frente a $0{,}76$) y todas quedan por debajo del SPY ($0{,}90$) en el
valor puntual, los contrastes de Memmel no permiten rechazar que tengan el mismo
Sharpe. Es coherente con la literatura \cite{DeMiguel2009} y tiene una
explicación clara: con 36 meses de ventana, el error de estimación de los 495
parámetros del modelo domina cualquier ganancia teórica de la optimización. La
propia $1/N$ lo ilustra: sin estimar ni optimizar nada, y con cero turnover,
logra un $10{,}7\%$ anualizado con un $14{,}1\%$ de volatilidad.

\subsection{Conclusión final}

La integración de Machine Learning en el modelo de Markowitz produce carteras
competitivas (la mejor, Lasso, iguala a la $1/N$) cuando la predicción es
agrupada y fuertemente regularizada, pero ninguna diferencia alcanza
significación estadística: el error de estimación sigue siendo el factor
limitante. Así, la principal contribución no es un modelo que gane al mercado,
sino mostrar con rigor \textit{por qué es tan difícil ganarle}: con datos
realistas, las diferencias entre estrategias son indistinguibles del ruido. Y
también qué técnicas ayudan a mitigar ese error de estimación, aunque no a
eliminarlo.

%--------------------------------------------------------------------
\section{Líneas futuras de investigación}
%--------------------------------------------------------------------

Los resultados de este trabajo sugieren varias direcciones de investigación:

\begin{itemize}
    \item \textbf{Ventanas de estimación más largas} (5--10 años), con un
          $n/T$ más favorable, a costa de mezclar regímenes de mercado.
    \item \textbf{Más señal en $\boldsymbol{\mu}$:} features fundamentales y
          macro (ratios contables, analistas, tipos, VIX) y técnicas de ML más
          potentes (Gradient Boosting, redes, LSTM), limitadas aquí por lo
          escaso de los datos (180 meses).
    \item \textbf{Tratar la incertidumbre de $\boldsymbol{\mu}$:} optimización
          robusta (minimizar el peor caso en una región de confianza) o
          enfoques bayesianos como Black-Litterman.
    \item \textbf{Covarianza dinámica} (GARCH multivariante, EWMA) para
          correlaciones cambiantes en regímenes de estrés.
    \item \textbf{Gestión del turnover:} penalizarlo en la función objetivo o
          como restricción, para limitar los costes de transacción.
    \item \textbf{Validación en otros mercados} y clases de activos, para medir
          la generalidad de los resultados.
\end{itemize}

En definitiva, optimizar carteras con parámetros estimados (no conocidos) sigue
siendo un problema abierto, en la frontera entre las matemáticas, la estadística
y las finanzas. Este trabajo ha querido aportar a ese campo rigor matemático,
honestidad empírica y una implementación reproducible.
